% !TEX root = ../main.tex

% 结论
\begin{conclusions}
随着推荐系统在各类互联网平台中的广泛应用，用户交互数据的多样化与复杂性不断提升，传统的单行为建模方式已难以充分刻画用户的真实偏好特征。多行为与多模态信息的引入为推荐系统带来了新的机遇，也带来了噪声干扰、语义偏移及鲁棒建模等新的挑战。本文围绕多行为与多模态推荐场景下的噪声抑制与偏好建模问题展开研究，提出了两种创新性的去噪推荐框架。

首先，提出了一种新颖的多行为推荐框架——NIPER，用于解决多行为推荐系统中噪声交互与真实偏好提取的核心问题。与以往方法普遍将所有辅助行为数据直接融合到目标行为模型中不同，NIPER能够区分辅助行为中的噪声与真实偏好。其中，行为间与行为内加权去噪模块通过为多行为交互分配权重，有效削弱了噪声交互的影响，确保无关或偶然的行为被过滤；基于偏好的表征增强模块则利用自监督学习机制，从辅助行为中提取真实偏好信息，从而增强模型的表达能力。广泛的实验结果表明，NIPER在推荐精度与鲁棒性方面均显著优于现有最先进方法，验证了该框架在多行为推荐任务中的有效性。

接着，本文针对多模态多行为推荐场景中的噪声干扰与鲁棒表示建模问题，提出了一种融合大语言模型（LLM）语义理解能力与图神经网络（GNN）结构建模优势的统一去噪推荐框架。该方法通过多模态内容辅助、语义引导消息传递与跨图对比学习三方面的设计，实现了多源信号下用户真实偏好的精准建模。首先，在多模态辅助偏好建模阶段，引入多模态大语言模型生成物品的语义摘要文本，并结合多行为交互生成用户偏好摘要，从而在语义层面实现跨模态对齐；随后，通过模态语义嵌入计算用户与物品相似度分布，设计硬去噪机制以识别并剔除低相关交互，显著提升交互图的纯净度。其次，在软去噪模块中，构建语义感知注意力机制，引导GNN在消息传递过程中自适应地聚合语义相关性更高的邻居节点，从而降低噪声扩散并增强表示的鲁棒性。最后，通过模态感知图增强与跨图对比学习模块，引入语义增强图以补充协同图中的偏好缺失信息，并在跨图约束下实现多视角表示一致性。实验结果表明，该模型在多个指标上均显著优于主流基线模型，验证了所提框架在融合多源模态信息、抑制交互噪声与提升推荐精度方面的优势。

研究仍然存在一些不足。本文提出的多行为与多模态去噪推荐方法在泛化性与可扩展性方面仍有进一步提升空间。首先，模型整体结构较为复杂，包含多层去噪与语义对齐模块，虽然能够有效提升推荐性能，但也带来了较高的计算开销与参数规模，在大规模推荐场景下的训练效率与推理速度仍需优化。最后，本文的实验主要基于公开数据集，缺乏在真实工业环境中的在线验证与用户反馈分析。未来的研究将探索基于因果推理的噪声建模与多模态语义对齐机制，结合任务自适应与轻量化建模策略，进一步提升模型的泛化性、可扩展性与稳健性，为多行为多模态推荐系统的发展提供更加通用且高效的解决方案。

\end{conclusions}
